\chapter{Examples of phase density and perimeter blow-ups}\label{sec:examples-perimeter-blowups}
These elementary examples separate the reduced boundary, phase density, topological boundary and perimeter tangent measures. They use the distributional framework of \cite{Sim18}, Chapter 3, Section 4, and the explicit blow-up results developed in the perimeter chapter. Their calculations do not use the general rectifiability theorem.

\begin{example}[A smooth subgraph has the expected oriented blow-up]\label{ex:smooth-subgraph-blowup}
Suppose $u\in C^1$ near $0\in\mathbb R^{n-1}$, $u(0)=0$, $Du(0)=0$, and locally $E=\{(z,s):s<u(z)\}$. The origin belongs to $\partial^*E$, its outward normal is $e_n$, and $E_{0,r}\to\{s<0\}$ locally in volume. Its rescaled perimeter measures converge to $\mathcal H^{n-1}\llcorner\{s=0\}$.
If $u$ is $C^2$, then for each fixed $R$,
\[
 |(E_{0,r}\mathbin{\triangle}\{s<0\})\cap B_R|=O_R(r),
 \qquad
 \frac{1}{|D\chi_E|(B_r)}\int_{B_r}|\nu-e_n|^2\,d|D\chi_E|=O(r^2).
\]
The constants depend on a local bound for $D^2u$.
\end{example}
\begin{proof}
The local smooth-domain calculation in Example \ref{ex:smooth-perimeter}, or direct integration by parts in the subgraph, gives outward normal $(-Du,1)/\sqrt{1+|Du|^2}$ and perimeter density $\sqrt{1+|Du|^2}\,dz$ on its graph. This continuous normal tends to $e_n$ and the graph meets every ball in positive surface measure, so its normalized polar average tends to $e_n$, proving reduced-boundary membership.

The dilated graph has height $u(rz)/r$ and gradient $Du(rz)$, both tending uniformly to zero on every fixed compact base set. The symmetric difference in $B_R$ is bounded by
\[
 \int_{|z|<R}\frac{|u(rz)|}{r}\,dz\longrightarrow0.
\]
For any compactly supported continuous scalar test $\phi$, its rescaled perimeter integral is
\[
 \int\phi(z,u(rz)/r)\sqrt{1+|Du(rz)|^2}\,dz,
\]
with integration confined to a fixed bounded projection of the test support. Dominated convergence gives the planar limit. When $|D^2u|\leq C$, Taylor's formula bounds $|u(rz)|/r$ by $Cr|z|^2/2$ and $|Du(z)|$ by $C|z|$. The first bound gives the stated volume rate. On the graph inside $B_r$, $|\nu-e_n|\leq C' r$; averaging its square gives the second rate. These are local rates for a smooth model, not uniform rates for an arbitrary finite-perimeter set.
\end{proof}

\begin{example}[A planar sector has unit perimeter density at every vertex angle]\label{ex:sector-reduced-boundary}
For $0<\theta<2\pi$, let
\[
 E_\theta=\{\rho(\cos t,\sin t):\rho>0,\ -\theta/2<t<\theta/2\}\subset\mathbb R^2.
\]
At the vertex,
\[
 |E_\theta\cap B_r|=\tfrac12\theta r^2,\qquad
 P(E_\theta;B_r)=2r,\qquad
 \frac{D\chi_{E_\theta}(B_r)}{P(E_\theta;B_r)}
 =\sin(\theta/2)e_1.
\]
Consequently the perimeter density normalized by $\omega_1r=2r$ is one for every angle, but the vertex is reduced-boundary exactly when $\theta=\pi$. Perimeter density one alone does not characterize the reduced boundary.
\end{example}
\begin{proof}
Polar integration gives the volume. On the upper ray the outward normal is $(-\sin(\theta/2),\cos(\theta/2))$; on the lower ray it is $(-\sin(\theta/2),-\cos(\theta/2))$. Integration by parts on the sector cut off between radii $\varepsilon$ and a large radius gives the derivative measure as minus these normals times length on the two rays. The test field is zero near the large boundary, and the small circular boundary term tends to zero with its length $\theta\varepsilon$. There is no vertex atom. Each ray has length $r$ inside $B_r$, so its variation mass is $2r$ and its vector derivative has value $2r\sin(\theta/2)e_1$. The normalized vector has length one precisely at $\theta=\pi$. At every other angle the limiting average exists but has length less than one, so Definition \ref{def:reduced-boundary} excludes the vertex.
\end{proof}

\begin{example}[Density one half does not characterize the reduced boundary]\label{ex:checkerboard-density-half}
For $E=\{(x,y):xy>0\}$, the origin has phase density $1/2$ and lies in the support of the perimeter measure, but does not belong to $\partial^*E$. In fact
\[
 P(E;B_r)=4r,\qquad D\chi_E(B_r)=0
\]
for every $r>0$.
\end{example}
\begin{proof}
The two occupied quadrants have total angle $\pi$, so they occupy half of every centered ball. Applying the ray calculation in Example \ref{ex:sector-reduced-boundary} separately to these two disjoint quadrants, or integrating by parts in each, gives one copy of length measure on each of the four axis rays. On opposite rays the outward normals are opposite, so the four vector contributions cancel in every centered ball. The positive variation is $4r$, whereas the averaged derivative vector is zero. This violates the unit-length condition in the reduced-boundary definition despite the exact phase density. The example is distinct from a smooth crossing immersion: it is a measurable phase whose distributional boundary has four rays.
\end{proof}

\begin{example}[A thin slab loses all phase volume but retains two sheets of perimeter]\label{ex:collapsing-slab-perimeter}
For $E_j=\{y\in\mathbb R^n:0<y_n<1/j\}$, one has locally
\[
 \chi_{E_j}\to0\text{ in }L^1,\qquad
 D\chi_{E_j}\rightharpoonup0,\qquad
 |D\chi_{E_j}|\rightharpoonup2\mathcal H^{n-1}\llcorner\{y_n=0\}.
\]
Thus local $L^1$ convergence and locally bounded perimeter do not ensure convergence of perimeter measures to the limit set's perimeter.
\end{example}
\begin{proof}
The volume on a fixed compact set is at most its bounded projection volume times $1/j$. By one-dimensional integration in $y_n$,
\[
 D\chi_{E_j}=e_n\mathcal H^{n-1}\llcorner\{y_n=0\}
             -e_n\mathcal H^{n-1}\llcorner\{y_n=1/j\}.
\]
The two planes are disjoint, so their positive measures add in the total variation. Continuous compactly supported test integrals on the translated plane converge to those on the original plane by uniform continuity and bounded projection. Consequently the vector measures cancel in the limit while their variations add. The missing hypothesis of Proposition \ref{prop:aligned-vector-variation-limit} is alignment of the two opposite polar directions.
\end{proof}

\begin{example}[A quadratic cusp has a double-line perimeter tangent and an empty phase tangent]\label{ex:quadratic-cusp-perimeter}
Let $E=\{(x,y)\in\mathbb R^2:|y|<x^2\}$. At the origin,
\[
 \chi_{E_{0,r}}\to0\quad\text{in }L^1_{\mathrm{loc}},\qquad
 r^{-1}(F_r^{-1})_\#|D\chi_E|
 \rightharpoonup2\mathcal H^1\llcorner\{y=0\}.
\]
The origin lies in the perimeter support and has phase density zero; it is not in the reduced boundary. In fact $D\chi_E(B_r)=0$ for every $r>0$. This is a phase-boundary example, distinct from the calibrated complex cusp in Example \ref{ex:complex-cusp-current}.
\end{example}
\begin{proof}
For a compactly supported smooth field $X=(X_1,X_2)$, integrate first in $y$ over $(-x^2,x^2)$ and use the one-dimensional Leibniz rule in $x$. One obtains
\[
 \int_E\operatorname{div}X
 =\int_{\mathbb R}\bigl[-2xX_1(x,x^2)+X_2(x,x^2)
                       -2xX_1(x,-x^2)-X_2(x,-x^2)\bigr]\,dx.
\]
Thus the perimeter is arclength on the two parabolas, with outward normals $(-2x,1)/\sqrt{1+4x^2}$ above and $(-2x,-1)/\sqrt{1+4x^2}$ below. They meet only at one point, which has zero length measure. Symmetry in $x$ cancels the horizontal vector integral in a centered ball, and symmetry in $y$ cancels the vertical integral. Their positive lengths do not cancel and are positive in every centered ball.

The rescaled phase is $\{(X,Y):|Y|<rX^2\}$, whose volume in $B_R$ is bounded by $2r\int_{-R}^R X^2\,dX\to0$. For a compactly supported continuous test $\phi$, its rescaled perimeter integral is
\[
 \int_{\mathbb R}\bigl[\phi(X,rX^2)+\phi(X,-rX^2)\bigr]
                       \sqrt{1+4r^2X^2}\,dX.
\]
Dominated convergence on a bounded projection gives twice line measure. The empty phase tangent and double-line perimeter tangent demonstrate precisely the cancellation that polar alignment excludes at a reduced-boundary point.
\end{proof}

\begin{example}[A negligible dense modification preserves every perimeter blow-up]\label{ex:halfspace-dense-modification}
Let $H=\{y_n<0\}$, and set
\[
 E=(H\setminus\mathbb Q^n)\cup(\mathbb Q^n\setminus H).
\]
Then $E=H$ almost everywhere, $\partial^*E=\{y_n=0\}$ and every reduced-boundary blow-up is the same oriented half-space as for $H$. Nevertheless the topological boundary of this representative is all of $\mathbb R^n$.
\end{example}
\begin{proof}
The symmetric difference is contained in the countable null set $\mathbb Q^n$. Characteristic functions, derivative measures and their polar averages are unchanged, so Lemma \ref{lem:halfspace-derivative-measure} gives the reduced boundary and its outward normal. At a point off the plane, some ball has zero perimeter measure, excluding it from the reduced boundary; on the plane the polar vector is constant of unit length. The set and its complement are both dense: every ball contains a smaller ball lying strictly in one of the two half-spaces, and within that smaller ball there are both rational and nonrational points, one type in $E$ and the other outside it. Hence every point is a topological boundary point. The result extends Example \ref{ex:null-modification-boundary} to a nonzero perimeter measure with unchanged geometric blow-ups.
\end{proof}
